\documentclass[10.5pt,a4paper]{article}

\usepackage[utf8]{inputenc}
\usepackage[indonesian]{babel}
\usepackage[a4paper,top=1.8cm,bottom=1.8cm,left=1.9cm,right=1.9cm]{geometry}
\usepackage{amsmath,amssymb}
\usepackage{graphicx}
\usepackage{booktabs}
\usepackage{array}
\usepackage{enumitem}
\usepackage{xcolor}
\usepackage{colortbl}
\usepackage{tcolorbox}
\usepackage{fancyhdr}
\usepackage{titlesec}

\definecolor{primaryblue}{RGB}{24, 43, 73}
\definecolor{secondaryblue}{RGB}{41, 128, 185}
\definecolor{boxbg}{RGB}{245, 248, 252}
\definecolor{goldaccent}{RGB}{211, 84, 0}
\definecolor{headerrow}{RGB}{24, 43, 73}

\pagestyle{fancy}
\fancyhf{}
\fancyhead[L]{\small\textbf{\color{primaryblue}Instrumen Umpan Balik Mahasiswa (OBE-CQI)}}
\fancyhead[R]{\small\color{gray}TEE-412 / Teknik Elektro FT UNIB}
\fancyfoot[L]{\footnotesize\color{gray}Evaluasi Penjaminan Mutu Berkelanjutan (LAM Teknik / IABEE)}
\fancyfoot[R]{\footnotesize Halaman \thepage}
\renewcommand{\headrulewidth}{0.5pt}
\renewcommand{\footrulewidth}{0.4pt}

\titleformat{\section}{\color{primaryblue}\normalfont\large\bfseries}{\thesection.}{0.5em}{}[\titlerule]
\setlist{nosep}
\newcommand{\hdr}[1]{\cellcolor{headerrow}\color{white}\textbf{#1}}
\newcommand{\likert}{$\square\,1 \quad \square\,2 \quad \square\,3 \quad \square\,4 \quad \square\,5$}

\begin{document}

\begin{center}
    \includegraphics[width=1.9cm]{assets/logo_unib.png}\\[0.2cm]
    {\large \textbf{\color{primaryblue}UNIVERSITAS BENGKULU -- FAKULTAS TEKNIK -- S1 TEKNIK ELEKTRO}}\\[0.15cm]
    {\LARGE \textbf{\color{primaryblue}KUESIONER UMPAN BALIK PEMBELAJARAN (OBE-CQI)}}\\[0.08cm]
    {\large \textbf{\color{goldaccent}ANALISIS DERET WAKTU DAN PERAMALAN (TEE-412)}}\\[0.05cm]
    {\small Evaluasi Tengah Semester (Minggu 8) \& Akhir Semester (Minggu 16) -- Tanpa Nama (\textit{Anonim})}
\end{center}

\vspace{0.15cm}
\begin{tcolorbox}[colback=boxbg,colframe=primaryblue,arc=2mm,title=\textbf{Petunjuk Pengisian \& Skala Penilaian (Likert 1--5)}]
\small
Kuesioner ini bersifat \textbf{anonim} dan \textbf{tidak memengaruhi nilai akhir}. Masukan Anda digunakan langsung dalam siklus \textit{Continuous Quality Improvement} (CQI / PDCA) untuk meningkatkan kualitas bahan ajar, praktikum Google Colab, Virtual Lab Studio, dan proyek akhir pada semester berjalan maupun berikutnya.\\[0.1cm]
\textbf{Periode Evaluasi:} $\square$ Tengah Semester (Minggu 8 -- CPMK 1--2) \qquad $\square$ Akhir Semester (Minggu 16 -- CPMK 3--5)\\
\textbf{Skala:} \textbf{1} = Sangat Tidak Setuju/Kurang \;|\; \textbf{2} = Tidak Setuju \;|\; \textbf{3} = Cukup \;|\; \textbf{4} = Setuju \;|\; \textbf{5} = Sangat Setuju/Sangat Baik.
\end{tcolorbox}

\section{Dimensi A: Persepsi Pencapaian CPMK (\textit{Indirect Assessment} OBE)}
{\small
\renewcommand{\arraystretch}{1.28}
\begin{tabular}{|p{0.6cm}|p{10.5cm}|p{4.8cm}|}
\hline
\hdr{No} & \hdr{Pernyataan Penguasaan Capaian Pembelajaran} & \hdr{Skala (1--5)} \\ \hline
A1 & Saya mampu mengidentifikasi komponen deret waktu (tren, musiman, residu), menguji stasioneritas (ADF/KPSS), dan membaca ACF/PACF (\textbf{CPMK-1}). & \likert \\ \hline
A2 & Saya mampu memodelkan peramalan statistik klasik (Holt-Winters, SARIMA/SARIMAX) serta menerapkan \textit{Discrete Kalman Filter} (\textbf{CPMK-2}). & \likert \\ \hline
A3 & Saya mampu menganalisis harmonisa dan sinyal transien kelistrikan pada domain frekuensi (FFT/PSD) dan waktu-frekuensi (Wavelet DWT) (\textbf{CPMK-3}). & \likert \\ \hline
A4 & Saya mampu membangun pipeline ML (XGBoost) dan DL (LSTM/1D-CNN/Autoencoder) bebas \textit{data leakage} untuk beban, PLTS, dan anomali (\textbf{CPMK-4}). & \likert \\ \hline
A5 & Saya mampu mengintegrasikan model peramalan pada studi kasus sistem tenaga riil dalam kerja tim serta menyusun laporan format IEEE (\textbf{CPMK-5}). & \likert \\ \hline
\end{tabular}
}

\section{Dimensi B: Kualitas Bahan Ajar, Kuis Self-Check, \& Portal Web}
{\small
\renewcommand{\arraystretch}{1.28}
\begin{tabular}{|p{0.6cm}|p{10.5cm}|p{4.8cm}|}
\hline
\hdr{No} & \hdr{Pernyataan Kualitas Ekosistem Bahan Ajar} & \hdr{Skala (1--5)} \\ \hline
B1 & Slide Beamer, Problem Sets (C1--C6), dan Worksheets tersusun sistematis serta relevan dengan aplikasi Teknik Elektro modern. & \likert \\ \hline
B2 & Kuis \textit{Self-Check} interaktif di portal web membantu saya mengecek pemahaman konsep sebelum praktikum maupun ujian. & \likert \\ \hline
B3 & Portal web (\texttt{ndaratha.my.id/deret-waktu}) mudah diakses, cepat dimuat, dan seluruh dokumen RPS/kisi-kisi mudah ditemukan. & \likert \\ \hline
\end{tabular}
}

\section{Dimensi C: Efektivitas Praktikum Google Colab \& Virtual Forecasting Studio}
{\small
\renewcommand{\arraystretch}{1.28}
\begin{tabular}{|p{0.6cm}|p{10.5cm}|p{4.8cm}|}
\hline
\hdr{No} & \hdr{Pernyataan Praktikum \& Lingkungan Komputasi} & \hdr{Skala (1--5)} \\ \hline
C1 & 12 notebook praktikum dapat dibuka 1-klik di Google Colab dan dijalankan lancar tanpa kendala instalasi/error dependensi. & \likert \\ \hline
C2 & Tiga dataset CSV terkurasi (beban regional, PLTS 500 kWp, smart meter) realistis dan membantu memahami perilaku sistem kelistrikan. & \likert \\ \hline
C3 & Virtual Forecasting Studio (termasuk fitur \textit{Upload CSV}) membantu memvisualisasikan pengaruh parameter ($Q, R, \alpha, \beta, \gamma$) secara intuitif. & \likert \\ \hline
\end{tabular}
}

\section{Dimensi D: Asesmen, Proyek Akhir IEEE, \& Umpan Balik Terbuka}
{\small
\renewcommand{\arraystretch}{1.28}
\begin{tabular}{|p{0.6cm}|p{10.5cm}|p{4.8cm}|}
\hline
\hdr{No} & \hdr{Pernyataan Asesmen \& Beban Belajar} & \hdr{Skala (1--5)} \\ \hline
D1 & Proporsi bobot penilaian (Kuis 10\%, Praktikum 25\%, UTS 20\%, Proyek 25\%, UAS 20\%) adil dan rubrik proyek IEEE transparan. & \likert \\ \hline
D2 & Beban kerja mingguan (kuliah, praktikum, dan proyek akhir tim) sesuai dengan bobot 3 SKS dan dapat diselesaikan tepat waktu. & \likert \\ \hline
\end{tabular}
}

\vspace{0.25cm}
\noindent\textbf{Pertanyaan Reflektif Terbuka (Isi singkat untuk perbaikan CQI):}
\begin{enumerate}[label=\textbf{\arabic*.}, leftmargin=1.5em, itemsep=6pt]
    \item Materi, rumus, atau bab praktikum mana yang paling sulit dipahami dan memerlukan contoh/pendalaman tambahan?\\
    \rule{\textwidth}{0.4pt}\\[0.25cm]
    \rule{\textwidth}{0.4pt}
    \item Kendala teknis apa (jika ada) yang Anda alami saat menjalankan notebook Google Colab, Virtual Studio, atau menyusun laporan LaTeX IEEE?\\
    \rule{\textwidth}{0.4pt}\\[0.25cm]
    \rule{\textwidth}{0.4pt}
    \item Fitur, dataset, atau topik rekayasa elektro apa yang paling bermanfaat dan/atau Anda sarankan untuk ditambahkan pada semester berikutnya?\\
    \rule{\textwidth}{0.4pt}\\[0.25cm]
    \rule{\textwidth}{0.4pt}
\end{enumerate}

\end{document}
